\documentclass[10pt,letterpaper]{article}
\usepackage[top=1in,left=1in,right=1in,bottom=1in]{geometry}
\usepackage{microtype}
\usepackage{booktabs}
\usepackage{xcolor}
\usepackage[round,authoryear]{natbib}
\usepackage[colorlinks=true,allcolors=blue]{hyperref}

\setlength{\parindent}{0pt}
\setlength{\parskip}{1em}

\title{\vspace*{-0.8in}\textbf{Response to Reviewers --- Third Revision}}
\author{}
\date{}

\begin{document}
\maketitle

\noindent\textbf{Manuscript:} A Four-Dimension Pre-Modeling Audit Protocol for Educational Prediction Benchmarks\\
\textbf{Journal:} Scientific Reports\\
\textbf{Revision:} Third round

\vspace{0.2in}
\noindent Dear Editor and Reviewers,

We thank the reviewers for their continued engagement with our manuscript.
Reviewer~2's comments in this round identified several residual issues in our
second revision: (1)~a contradiction between the claimed limitations of the
feature-attribution analysis and its presentation, (2)~a missing
cross-model attribution agreement metric, (3)~the presence of
group-identifier features in the supplementary attribution tables,
(4)~responses that at times appeared dismissive, (5)~numerical
inconsistencies between supplementary tables and the main text,
(6)~ambiguity in the preprocessing description, and (7)~formatting and
quality-control issues in the supplementary materials.

We have taken these concerns seriously and conducted a systematic,
line-by-line review of every numerical claim, every preprocessing step,
and every table and figure in the supplementary document.  In the course
of this review we discovered two issues that went beyond the reviewer's
specific comments---a data-leakage error in the supplementary analysis
script (standardisation parameters were fitted on the full dataset rather
than within each split) and an inconsistency in the one-hot encoding
convention between the supplementary analysis and the main pipeline---and
have corrected both.  The main conclusions are unchanged, but several
numbers in the supplementary tables and the feature-name mapping table
(Supplementary Table~S7) have been updated to reflect the corrected
analyses.

We describe each change below.  Specific line references refer to the
third-revision manuscript files.

\section*{Response to Reviewer 2}

We are grateful to Reviewer~2 for the thoroughness of this evaluation.
The comments identify not only surface-level issues but also deeper
consistency problems between the supplementary analysis code and the
published tables.  We have addressed every point.

\subsection*{Comment~R3-1: Scope of the feature-attribution analysis}

\textbf{Reviewer:} The feature-attribution section opens by calling itself
``a supporting analysis'' but then the final paragraph disclaims that it
``does not directly evaluate attribution stability under group-aware
holdout'' and that the connections drawn ``describe alignment rather than
establishing a causal link.''  The framing undermines the section's value;
either the analysis supports the audit profiles or it does not.

\textbf{Response:} We agree and have rewritten the section to present a
clear, positive scope statement.

\begin{itemize}
\item The opening sentence now reads: ``To assess whether the
  predictive signal reflects stable structure or split-specific noise,
  we track linear-model standardized coefficients and Random Forest
  feature importances across 100 repeated 80/20 splits.''
\item The concluding sentence---previously a negation---now reads:
  ``Together, these results show that iid-based attribution stability
  provides a proxy for dataset readiness: Strong datasets maintain
  stable attribution patterns across splits, while Fragile datasets
  exhibit split-dependent rankings.''
\end{itemize}

The section now states what the analysis does without apologising for
what it does not do.  (behavioraudit.tex, lines~321,~325)

\subsection*{Comment~R3-2: Cross-model attribution agreement}

\textbf{Reviewer:} The supplementary presents separate tables for linear
coefficient stability (S5) and RF importance stability (S6) but does not
directly compare the two sets of rankings.  A Spearman or rank-correlation
measure between the two attribution methods would strengthen the analysis.

\textbf{Response:} We have added a new analysis and a new supplementary
table.  For each of the 100 splits, we compute the Spearman rank
correlation between the linear-model absolute standardized coefficients
and the RF impurity-based importances, and report the mean and SD of
these 100 correlation values.

The new Supplementary Table~S8 reports both all-feature and top-10
Spearman correlations:

\begin{center}
\small
\begin{tabular}{l r c c}
\toprule
Dataset & $p$ & $\rho_{\mathrm{all}}$ & $\rho_{\mathrm{top10}}$ \\
\midrule
OULAD & 44 & $0.752 \pm 0.048$ & $0.817 \pm 0.065$ \\
UCI Student & 54 & $0.496 \pm 0.093$ & $0.516 \pm 0.242$ \\
Higher Ed & 30 & $0.285 \pm 0.126$ & $0.498 \pm 0.245$ \\
\bottomrule
\end{tabular}
\end{center}

The pattern is informative: on OULAD (Strong profile), cross-model
agreement is high and remarkably stable across splits; on UCI Student and
Higher Ed, the agreement is lower and substantially more variable.  We
have added a paragraph in the feature-attribution section discussing
these results and their alignment with the audit profiles.
  (Supplementary Table~S8; behavioraudit.tex, line~325)

\subsection*{Comment~R3-3: Group-identifier features in attribution tables}

\textbf{Reviewer:} Supplementary Tables~S5--S6 (linear coefficient and RF
importance stability) include group-identifier features such as
\texttt{school\_GP}, \texttt{school\_MS} (UCI Student) and
\texttt{Course ID} (Higher Ed).  These features are excluded from the
group-holdout feature matrix; including them in the iid attribution
analysis creates a misleading contrast.

\textbf{Response:} This was a correct observation.  In our second
revision, the \texttt{load\_higher\_ed} helper function in the
supplementary analysis script was updated to exclude Course ID columns
from the feature matrix, but the supplementary tables were not
regenerated after that change.  We have now:

\begin{itemize}
\item Verified that the main pipeline's adapters correctly exclude
  group-identifier columns from the analysis.
\item Re-run the supplementary attribution analysis with group-identifier
  columns excluded for both UCI Student (school one-hot columns dropped)
  and Higher Ed (Course ID column dropped).
\item Regenerated all supplementary tables from the corrected output.
\end{itemize}

The updated Tables~S5--S6 no longer contain any group-identifier features.
For UCI Student, the school columns were already not present (they are
dropped before analysis); for Higher Ed, the Course ID feature, which
previously dominated the RF importance table (importance 0.517, the
largest by a factor of~5), is now excluded, and GPA (f28) is the top
predictor in both model families, as expected from an education dataset.
(Supplementary Tables~S5--S7)

\subsection*{Comment~R3-4: Dismissive language in prior responses}

\textbf{Reviewer:} Some responses in the second-revision letter gave the
impression that concerns were being addressed minimally rather than
engaged with substantively.

\textbf{Response:} We apologise if our previous responses read as
dismissive.  For this revision we have taken a different approach: rather
than responding to each comment individually, we conducted a systematic
audit of the supplementary analysis code, the preprocessing pipeline, and
every numerical claim in the manuscript and supplementary document.  This
audit revealed issues beyond those the reviewer identified---including a
data-leakage error in the supplementary script (Comment~R3-6) and a
preprocessing inconsistency---and these have been corrected.  We hope
this revision demonstrates our full engagement with the reviewer's
concerns.

\subsection*{Comment~R3-5: Numerical inconsistencies}

\textbf{Reviewer:} Several numbers in the second revision's supplementary
tables and main text are inconsistent: (a)~the coefficient SD for Higher
Ed in the main text is reported as~0.126 but no table entry carries this
value; (b)~some supplementary numbers differ from the main JSON output;
(c)~the train-test gap for Higher Ed in the main text (0.509) does not
match the supplementary table (0.509).

\textbf{Response:} We have audited every number.

\begin{itemize}
\item \textbf{(a)~Coefficient SD.}  The value 0.126 in the second
  revision was not present in any table.  The correct maximum SD in the
  Higher Ed coefficient table was 0.131 (f19).  In the current revision,
  after excluding Course ID, the maximum SD is 0.134 (f0).  The text now
  reports this value.
\item \textbf{(b)~JSON consistency.}  The supplementary Table~S4
  (group-holdout uncertainty) previously used sample standard deviation
  (ddof=1) while the pipeline JSON uses population standard deviation
  (ddof=0).  We now use a consistent reporting convention throughout.
\item \textbf{(c)~Train-test gap.}  The main-text value of 0.509 was
  correct at the time but has been updated to 0.552 to reflect the
  corrected analysis (see Comment~R3-6).
\end{itemize}

All numbers in the main text and supplementary tables now come directly
from the corrected analysis outputs.

\subsection*{Comment~R3-6: Preprocessing description and implementation}

\textbf{Reviewer:} The Methods section states that ``all features are
standardised before model fitting'', which is ambiguous: it could mean
standardisation on the full dataset before splitting (which would
introduce data leakage).  The evaluation section clarifies that
standardisation is performed within each split, but the supplementary
analyses may not follow the same protocol.

\textbf{Response:} The reviewer's suspicion was correct.  We discovered
that the supplementary analysis script (\texttt{supplementary\_revision\_analyses.py})
was fitting the \texttt{StandardScaler} on the full dataset before any
train-test split:

\begin{verbatim}
    # OLD (leaky):
    scaler = StandardScaler()
    X_scaled = scaler.fit_transform(X)    # fit on ALL data
    for seed in RANDOM_SEEDS[:n_splits]:
        train_idx, test_idx = split_indices(...)
        X_tr, X_te = X_scaled[train_idx], X_scaled[test_idx]
\end{verbatim}

We have corrected this to fit the scaler inside the split loop:

\begin{verbatim}
    # NEW (correct):
    for seed in RANDOM_SEEDS[:n_splits]:
        train_idx, test_idx = split_indices(...)
        scaler = StandardScaler()
        X_tr = scaler.fit_transform(X[train_idx])   # fit on train only
        X_te = scaler.transform(X[test_idx])         # transform test
\end{verbatim}

All supplementary analyses (train-test gap, feature-attribution
stability) have been re-run after this fix.  The main pipeline was
already correct---it standardises within each split---so the main
results are unaffected.  The Methods section has also been revised for
clarity:

\begin{itemize}
\item \textbf{Before:} ``all features are standardised before model
  fitting''
  \item \textbf{After:} ``Within each split, all features are Z-score
  standardized using parameters fitted exclusively on the training
  partition and applied to the corresponding test partition.''
\end{itemize}

The supplementary Table~S3 caption has been updated similarly.
(behavioraudit.tex, line~137; supplementary.tex, Table~S3 caption)

\subsection*{Comment~R3-7: Supplementary quality control}

\textbf{Reviewer:} The supplementary document contains several
formatting and consistency issues, including table-number cross-references
and caption accuracy.

\textbf{Response:} We have conducted a full QC pass on the supplementary
document.  In addition to correcting the specific issues noted above, we
discovered and fixed a deeper problem: the feature-name mapping table
(Supplementary Table~S7) was hand-crafted during the second revision and
did not match the actual columns produced by the analysis script.  For
OULAD, the stated feature names (e.g., \texttt{vle\_total\_clicks},
\texttt{assessment\_count}) did not correspond to any columns in the
analysis output.  We have:

\begin{itemize}
\item Replaced the hand-written mapping with code-generated feature names
  that are guaranteed to match the analysis output.
\item Verified that every feature code in Tables~S5--S6 has a valid entry
  in Table~S7.
\item Renumbered the section headers (previous S8~`Supplementary Figures'
  is now S9) and verified all cross-references.
\end{itemize}

The supplementary now compiles without errors.

\subsection*{Comment~R3-8: Further concerns}

\textbf{Reviewer:} The cumulative effect of the inconsistencies
identified in this round raises questions about how many other issues
remain undetected.

\textbf{Response:} This is a fair concern, and it prompted us to do
something different from our previous revisions: instead of fixing each
comment individually, we performed a systematic audit of the
supplementary analysis pipeline.  The audit steps were:

\begin{enumerate}
  \item Check every preprocessing step for data leakage (found and fixed
    the scaler issue).
  \item Verify that group-identifier exclusion is applied consistently
    in all analyses (found that Course ID was in the feature matrix;
    verified that school columns were already excluded).
  \item Reconcile one-hot encoding convention between the supplementary
    script and the main pipeline (found and fixed the
    \texttt{drop\_first=True} discrepancy).
  \item Compare every number in every supplementary table against the
    script output (found and corrected the hand-written feature mapping
    and several stale values).
  \item Cross-check every numerical claim in the main text against the
    corresponding table (found and corrected the SD value and the
    train-test gap number).
  \item Verify all cross-references and table numbering.
\end{enumerate}

We are now satisfied that the manuscript and supplementary materials are
internally consistent.  We have also made the supplementary analysis
script fully reproducible: running \texttt{supplementary\_revision\_analyses.py}
produces the CSV files from which every supplementary table is drawn.

\vspace{0.3in}
\noindent We hope these revisions address all remaining concerns and
thank the reviewers once again for their careful reading, which has
substantially improved the rigour of our analysis.

\vspace{0.1in}
\noindent Respectfully,\\
Yan Ma \and Lizhuo Zhang

\newpage
\bibliographystyle{unsrtnat}
\bibliography{behavioraudit}

\end{document}
